\begin{tcolorbox}[colback=red!3!white,colframe=red!50!black,boxrule=0.8pt,title={Korrekturhinweise},fonttitle=\bfseries]
\textbf{Korrekturen aus der Rechenprüfung (30.~Sept.~2026).} Die folgenden Stellen sind im Text korrigiert bzw. vermerkt; jede trägt dort einen datierten Hinweis mit dem vorherigen Wert:
\begin{itemize}
\item Interpretation der Tabelle: Die Drop-Werte ($2{,}935914$; $1{,}435758$; $51\,\%$) stimmen. Die bessere Robustheit gilt aber nur relativ -- absolut hat T0-QAT in beiden Spalten den größeren Fehler ($9{,}50$ statt $0{,}32$ in voller Präzision, $10{,}94$ statt $3{,}25$ quantisiert, Faktor $3{,}4$).
\item \enquote{Ohne empirisches Tuning}: Die Rauschamplitude ist $\xi \cdot 80000 = 10{,}67$ mit einem frei gewählten, datensatzspezifischen Skalierungsfaktor; $\xi$ legt die Amplitude nicht allein fest. $\xi$ selbst bleibt aus der Theorie hergeleitet.
\end{itemize}
\end{tcolorbox}

\section*{Abstract}
		Dieses Dokument stellt eine experimentelle Erprobung von $\xi$-bewusstem Quantization-aware Training vor, wobei $\xi = \frac{4}{3} \times 10^{-4}$ aus den physikalischen Grundprinzipien der T0-Theorie (Zeit-Masse-Dualität) hergeleitet ist. Die vorläufigen Ergebnisse zeigen einen relativ geringeren Leistungsabfall unter Quantisierungsrauschen als der Standardansatz, bei absolut höherem Fehler; die Rauschamplitude enthält zusätzlich einen datensatzspezifischen Skalierungsfaktor. \textit{(Korrigiert am 30.~Sept.~2026: vorher \enquote{verbesserte Robustheit gegenüber Standardansätzen} und \enquote{prinzipielle Rauschregularisierung} -- siehe Abschnitt Interpretation.)}

	
	
	\section{Einleitung}
	
	Quantization-aware training (QAT) hat sich als entscheidende Technik für das Deployment von neuronalen Netzen auf ressourcenbeschränkten Geräten etabliert. Allerdings basieren aktuelle Ansätze oft auf empirischen Rausch-Injektionsstrategien ohne theoretische Grundlage. Diese Arbeit führt $\xi$-aware QAT ein, basierend auf der T0 Zeit-Masse-Dualitätstheorie, die eine fundamentale physikalische Konstante $\xi$ bereitstellt, die numerische Präzisionsgrenzen natürlich regularisiert.
	
	\section{Theoretische Grundlagen}
	
	\subsection{T0 Zeit-Masse-Dualitätstheorie}
	
	Der Parameter $\xi = \frac{4}{3} \times 10^{-4}$ ist keine empirische Optimierung, sondern leitet sich aus ersten Prinzipien der T0-Theorie der Zeit-Masse-Dualität ab. Diese fundamentale Konstante repräsentiert den minimalen Rauschpegel, der physikalischen Systemen inhärent ist, und bietet eine natürliche Regularisierungsgrenze für numerische Präzisionslimits.
	
	Die vollständige theoretische Herleitung ist im T0 Theory GitHub Repository verfügbar\footnote{}, einschließlich:
	\begin{itemize}
		\item Mathematische Formulierung der Zeit-Masse-Dualität
		\item Herleitung fundamentaler Konstanten
		\item Physikalische Interpretation von $\xi$ als Quantenrauschgrenze
	\end{itemize}
	
	\subsection{Implikationen für AI Quantization}
	
	Im Kontext der Neural Network Quantization repräsentiert $\xi$ die fundamentale Präzisionsgrenze, unterhalb derer weitere Bit-Reduzierung aufgrund physikalischer Rauschbeschränkungen abnehmende Erträge liefert. Durch die Einbeziehung dieser physikalischen Konstante während des Trainings lernen Modelle, optimal innerhalb dieser natürlichen Präzisionsgrenzen zu operieren.
	
	\section{Experimenteller Aufbau}
	
	\subsection{Methodik}
	
	Wir entwickelten ein vergleichendes Framework zur Evaluierung von $\xi$-aware Training gegenüber standard Quantization-aware Ansätzen. Das experimentelle Design besteht aus:
	
	\begin{itemize}
		\item \textbf{Baseline:} Standard QAT mit empirischer Rausch-Injektion
		\item \textbf{T0-QAT:} $\xi$-aware Training mit physikalisch-informiertem Rauschen
		\item \textbf{Evaluation:} Quantisierungsrobustheit unter simulierter Präzisionsreduktion
	\end{itemize}
	
	\subsection{Datensatz und Architektur}
	
	Für die initiale Validierung verwendeten wir eine synthetische Regressionsaufgabe mit einer einfachen neuronalen Architektur:
	
	\begin{itemize}
		\item \textbf{Datensatz:} 1000 Samples, 10 Features, synthetisches Regressionsziel
		\item \textbf{Architektur:} Einzelne lineare Schicht mit Bias
		\item \textbf{Training:} 300 Epochen, Adam Optimizer, MSE Loss
	\end{itemize}
	
	\section{Ergebnisse und Analyse}
	
	\subsection{Quantitative Ergebnisse}
	
	\begin{table}[h]
		\centering
		\adjustbox{max width=\textwidth}{%
\begin{tabular}{lccc}
			\toprule
			\textbf{Methode} & \textbf{Volle Präzision} & \textbf{Quantisiert} & \textbf{Drop} \\
			\midrule
			Standard QAT & 0{,}318700 & 3{,}254614 & 2{,}935914 \\
			T0-QAT ($\xi$-aware) & 9{,}501066 & 10{,}936824 & 1{,}435758 \\
			\bottomrule
		\end{tabular}%
}
		\caption{Leistungsvergleich unter Quantisierungsrauschen}
		\label{tab:results}
	\end{table}
	
	\subsection{Interpretation}
	
	Die experimentellen Ergebnisse demonstrieren:
	
	\begin{itemize}
		\item \textbf{Relativ geringerer Drop:} T0-QAT zeigt eine um 51\,\% kleinere Leistungsverschlechterung durch die Quantisierung ($1{,}44$ statt $2{,}94$)
		\item \textbf{Absolut höherer Fehler:} In beiden Spalten liegt der Fehler von T0-QAT über dem Standardansatz ($9{,}50$ statt $0{,}32$ in voller Präzision, $10{,}94$ statt $3{,}25$ quantisiert, Faktor $3{,}4$)
		\item \textbf{Rauschresilienz:} Mit $\xi$-aware Rauschen trainierte Modelle reagieren weniger empfindlich auf Präzisionsvariationen in niedrigeren Bits
		\item \textbf{Rolle von $\xi$:} Die eingesetzte Rauschamplitude ist $\xi \cdot 80000 = 10{,}67$; der Skalierungsfaktor $80000$ ist datensatzspezifisch gewählt, $\xi$ legt die Amplitude also nicht allein fest
	\end{itemize}
	\textit{(Korrigiert am 30.~Sept.~2026: vorher \enquote{Verbesserte Robustheit: signifikant reduzierte Leistungsverschlechterung} ohne Hinweis auf den absolut höheren Fehler, und \enquote{effektive Regularisierung ohne empirisches Tuning} -- der Code verwendet \texttt{xi\_scaling = 80000.0}.)}
	
	\section{Implementierung}
	
	\subsection{Kernalgorithmus}
	
	Der T0-QAT Ansatz modifiziert Standard-Training durch Injektion von physikalisch-informiertem Rauschen während des Forward Pass:
	
	\begin{verbatim}
		# Fundamentale Konstante aus T0 Theorie
		xi = 4.0/3 * 1e-4
		
		def forward_with_xi_noise(model, x):
		weight = model.fc.weight
		bias = model.fc.bias
		
		# Physikalisch-informierte Rausch-Injektion
		noise_w = xi * xi_scaling * torch.randn_like(weight)
		noise_b = xi * xi_scaling * torch.randn_like(bias)
		
		noisy_w = weight + noise_w
		noisy_b = bias + noise_b
		
		return F.linear(x, noisy_w, noisy_b)
	\end{verbatim}
	
	\subsection{Vollständiger Experimenteller Code}
	
	\begin{verbatim}
		import torch
		import torch.nn as nn
		import torch.optim as optim
		import torch.nn.functional as F
		
		# xi aus T0-Theorie (Zeit-Masse-Dualität)
		xi = 4.0/3 * 1e-4
		
		class SimpleNet(nn.Module):
		def __init__(self):
		super().__init__()
		self.fc = nn.Linear(10, 1, bias=True)
		
		def forward(self, x, noisy_weight=None, noisy_bias=None):
		if noisy_weight is None:
		return self.fc(x)
		else:
		return F.linear(x, noisy_weight, noisy_bias)
		
		# T0-QAT Training Loop
		def train_t0_qat(model, x, y, epochs=300):
		optimizer = optim.Adam(model.parameters(), lr=0.005)
		xi_scaling = 80000.0  # Datensatz-spezifische Skalierung
		
		for epoch in range(epochs):
		optimizer.zero_grad()
		weight = model.fc.weight
		bias = model.fc.bias
		
		# Physikalisch-informierte Rausch-Injektion
		noise_w = xi * xi_scaling * torch.randn_like(weight)
		noise_b = xi * xi_scaling * torch.randn_like(bias)
		noisy_w = weight + noise_w
		noisy_b = bias + noise_b
		
		pred = model(x, noisy_w, noisy_b)
		loss = F.mse_loss(pred, y)
		loss.backward()
		optimizer.step()
		
		return model
	\end{verbatim}
	
	\section{Diskussion}
	
	\subsection{Theoretische Implikationen}
	
	Der relativ geringere Drop von T0-QAT ist ein Anlass zu prüfen, ob physikalische Prinzipien AI-Optimierungsstrategien informieren können; ein Beleg ist er bei absolut höherem Fehler und frei gewähltem Skalierungsfaktor noch nicht. Die $\xi$-Konstante könnte bieten:
	
	\begin{itemize}
		\item \textbf{Prinzipielle Regularisierung:} Physikalisch-basierte Alternative zu empirischen Methoden
		\item \textbf{Optimale Präzisionsgrenzen:} Natürliche Limits für Quantisierungs-Bit-Breiten
		\item \textbf{Cross-Domain Validierung:} Verbindung zwischen physikalischen Theorien und AI-Effizienz
	\end{itemize}
	
	\subsection{Praktische Anwendungen}
	
	\begin{itemize}
		\item \textbf{Low-Precision Inference:} INT4/INT3/INT2 Deployment (noch nicht getestet)
		\item \textbf{Edge AI:} Ressourcenbeschränktes Model Deployment
		\item \textbf{Quantum-Classical Interface:} Brückenschlag zwischen Quantenrauschmodellen und klassischer AI
	\end{itemize}
	
	\section*{Reproduzierbarkeit}
	
	Vollständiger Code, experimentelle Daten und theoretische Herleitungen sind in den assoziierten GitHub Repositories verfügbar:
	
	\begin{itemize}
		\item \textbf{Theoretische Grundlage:} 
	\end{itemize}
	
	\begin{thebibliography}{9}
		\bibitem{t0theory} 
		Pascher, J. \textit{T0 Time-Mass Duality Theory}. 
		GitHub Repository, 2025.
		
		\bibitem{qat} 
		Jacob, B. et al. \textit{Quantization and Training of Neural Networks for Efficient Integer-Arithmetic-Only Inference}. 
		CVPR, 2018.
		
		\bibitem{physicsai}
		Carleo, G. et al. \textit{Machine learning and the physical sciences}. 
		Reviews of Modern Physics, 2019.
	\end{thebibliography}
	
	\appendix
	\section{Theoretische Herleitungen}
	
	Vollständige mathematische Herleitungen der $\xi$-Konstante und T0 Zeit-Masse-Dualitätstheorie werden im dedizierten Repository gepflegt. Dies beinhaltet:
	
	\begin{itemize}
		\item Herleitung fundamentaler Gleichungen
		\item Konstanten-Berechnungen
		\item Physikalische Interpretationen
		\item Mathematische Beweise
	\end{itemize}
	